\subsubsection{Productivity}
\label{sec:semantics:grsync:theorems:productivity}

A productive component is one that is guaranteed to always produce an output for every possible
input it receives. In other words, it never gets stuck or fails to respond. This assumes that no
runtime errors occur, such as division by zero, overflows, or other exceptions.

\newcommand{\obsCompName}[1]{\ensuremath{\llbracket #1 \rrbracket^\textit{obs}_{\tyCtx, G}}}
\newcommand{\obsComp}[3]{\ensuremath{\obsCompName{#1}(#2, #3)}}

To formalize the notion of productivity, we first define the observation function for a component
$f$. The observation of a component $f$ at the $n^\text{th}$ time step on the input lists $L^+$ of
length at least $n$, denoted by \obsComp{f}{L^+}{n} and defined by the rule
\nameref{rule:semantics:grsync:theorems:productivity:observe:component}, is the execution of $n$
steps of the operational semantics of $f$ on the successive inputs of $L^+$, starting from its
initial state.

\begin{mathpar}
    \myinfer
    { s_0 = \initsync{f} \\ \forall k \in [0, n-1],\,\compOpsem{s_k}{(L[k])^+}{f}{s_{k+1}}{v^+_k} }
    { \obsComp{f}{L^+}{n} = v^+_{n-1} }
    {{\tiny Observe}Comp}
    \label{rule:semantics:grsync:theorems:productivity:observe:component}
\end{mathpar}

\Cref{thm:semantics:grsync:theorems:productivity} ensures that the operational semantics of a
well-defined component (\ie well-typed and causal) is productive by guaranteeing that its
observation function is well-defined.
%
The proof of this theorem proceeds by structural induction on the syntax of \GRSync. Due to the
presence of component calls within expressions, the proof is mutually recursive with the proofs
of the intermediate \Crefrange{lem:semantics:grsync:theorems:productivity:expression}
{lem:semantics:grsync:theorems:productivity:equations}, which establish equivalent productivity
properties for each syntactic construct of \GRSync (expressions, patterns, event patterns, and
equations). This mutual recursion is well-founded because, while expressions can contain
component calls, the components themselves are not mutually recursive, ensuring that the
induction remains valid.

\begin{lemma}[Productivity of \GRSync expression]
    \label{lem:semantics:grsync:theorems:productivity:expression}
    Let $e$ be a well-typed \GRSync expression, and let $\rho$ be an environment in which all
    dependencies of $e$ (\ie all identifiers in \deps{e}) have well-defined, non-$\bot$ values.
    %
    Assuming that no runtime errors occur, then for all state $s$ corresponding to $e$, the step
    function of $e$ in $\rho$ from $s$ produces a new compatible state $s'$ and a well-defined
    value.
\end{lemma}
\begin{skippableproof}
    We proceed by structural induction on expressions and rely on the induction hypothesis on
    component calls (recursive calls are forbidden).

    \paragraph{Constants}
    Let $c$ be a constant, let $s$ be a compatible state for the expression, let $\rho$ be an
    environment such that all dependencies of the expression are well-defined in $\rho$.
    %
    \Cref{eq:semantics:grsync:opsem:init:e:const} ensures that $s$ equals $()$.
    %
    The step function of the expression $c$ always produces the value $c$, as state by the rule
    \nameref{rule:semantics:grsync:opsem:step:e:const}, and $c$ is well-defined. Moreover, the produced
    state remains empty.

    \paragraph{Identifiers}
    Let $x$ be an identifier expression, let $s$ be a compatible state, let $\rho$ be an environment
    such that all dependencies of the expression are well-defined in $\rho$.
    %
    \Cref{eq:semantics:grsync:opsem:init:e:ident} ensures that $s$ equals $()$.
    %
    \Cref{rule:semantics:grsync:causality:dependencies:deps:e:ident} implies that $x \in \deps{x}$,
    so we know that $\rho[x]$ is assumed well-defined.
    %
    The rule \nameref{rule:semantics:grsync:opsem:step:e:ident} states that the value produced by
    the step function of $x$ is equal to $\rho[x]$, which is well-defined. Moreover, the produced
    state remains empty.

    \paragraph{Last memories}
    Let the expression \last{x} be the memory of the identifier $x$, let $s$ be a compatible state,
    let $\rho$ be an environment such that all dependencies of the expression are well-defined in
    $\rho$.
    %
    \Cref{eq:semantics:grsync:opsem:init:e:last} ensures that $s$ equals $()$.
    %
    \Cref{rule:semantics:grsync:causality:dependencies:deps:e:last} implies that $\last{x} \in
        \deps{\last{x}}$, so we know that $\rho[\last{x}]$ is assumed well-defined.
    %
    The rule \nameref{rule:semantics:grsync:opsem:step:e:last} states that the value produced by the
    step function of \last{x} is equal to $\rho[\last{x}]$, which is well-defined. Moreover, the produced
    state remains empty.

    \paragraph{Event emission}
    Let \emit{e} be the event emission of the expression $e$, let $s$ be a compatible state for
    the expression, let $\rho$ be an environment such that all dependencies of the expression are
    well-defined in $\rho$.
    %
    \Cref{eq:semantics:grsync:opsem:init:e:emit} ensures that $s$ is compatible with $e$.
    %
    \Cref{rule:semantics:grsync:causality:dependencies:deps:e:emit} ensures that all the
    dependencies of $e$ have well-defined values in the environment $\rho$.
    %
    By the induction hypothesis, evaluating $e$ with $\rho$ yields a well-defined value $v$ and a
    compatible state $s'$. The rule \nameref{rule:semantics:grsync:opsem:step:e:emit:value} implies
    that the value produced by the step function of \emit{e} is equal to \some{v}, which is
    well-defined. Moreover, the state produced remains compatible.

    \paragraph{Operations}
    Let $op(e^+)$ be the application of the operation $op$ on the expressions $e^+$, let $s$ be a
    compatible state for the expression, let $\rho$ be an environment such that all dependencies of
    the expression are well-defined in $\rho$.
    %
    \Cref{eq:semantics:grsync:opsem:init:e:op} ensures there exists a list of states $s_e^+$
    compatible with $e^+$ such that $s = (s_e^+)$.
    %
    \Cref{rule:semantics:grsync:causality:dependencies:deps:e:op} ensures that all the dependencies
    of the expressions $e^+$ have well-defined values in $\rho$.
    %
    By the induction hypothesis, each $e$ evaluates to a well-defined value $v$ and a new compatible
    state $s_e'$. The rule \nameref{rule:semantics:grsync:opsem:step:e:op:value} implies that the
    value produced by the step function of $op(e^+)$ is equal to $op(v^+)$, which is well-defined if
    and only if the expression $op(e^+)$ is well-typed and no runtime errors occur (which is the
    case by hypothesis). Moreover, the state produced remains compatible.

    \paragraph{Component calls}
    Let $f(e^+)$ be the call of the component $f$ on the expressions $e^+$, let $s$ be a compatible
    state for the expression, let $\rho$ be an environment such that all dependencies of the
    expression are well-defined in $\rho$.
    %
    \Cref{eq:semantics:grsync:opsem:init:e:app} ensures there exists a list of states $s_e^+$
    compatible with $e^+$ and a state $s_f$ compatible with $f$ such that $s = (s_f, s_e^+)$.
    %
    \Cref{rule:semantics:grsync:causality:dependencies:deps:e:app} ensures that all the dependencies
    of the expressions $e^+$ have well-defined values in $\rho$.
    %
    By the induction hypothesis, each $e$ evaluates to a well-defined value $v_i$ and a new
    compatible state $s_e'$. The induction on component calls ensures that the values produced by
    the step function of $f$ with the well-defined values $v_i^+$ yields the well-defined outputs
    $v_o^+$ and a compatible state $s_f'$. The rule
    \nameref{rule:semantics:grsync:opsem:step:e:app:value}
    implies that the values produced by the step function of $f(e^+)$ are equal to $v_o^+$, which
    are well-defined. Moreover, the state produced remains compatible.

    \setcounter{paragraph}{0}
\end{skippableproof}

\begin{lemma}[Productivity of \GRSync patterns]
    \label{lem:semantics:grsync:theorems:productivity:patterns}
    Let \guarded{pat^+ = v^+}{e_g} be a well-typed \GRSync pattern guarded by $e_g$ and matching the
    values $v^+$, let $\rho$ be an environment such that all dependencies of $e_g$ except those from
    the pattern $pat^+$ (\ie all identifiers in $\deps{e_g} \backslash \defs{pat^+}$) have
    well-defined values in the environment $\rho$ (non-$\bot$).
    %
    Assuming that no runtime errors occur, the step function of \guarded{pat^+ = v^+}{e_g} in $\rho$
    from an empty state produces another empty state, a boolean $b$, and an environment $\rho'$
    where identifiers defined by $pat^+$ (\ie all identifiers in \defs{pat^+}) are well-defined if
    $b$ is \true.
\end{lemma}
\begin{skippableproof}
    We proceed by structural induction on patterns.

    \paragraph{Guarded patterns}
    Let \guarded{pat^+ = v^+}{e_g} be a guarded tuple of patterns, let $\rho$ be an environment such
    that all dependencies of $e_g$ except those from the pattern $pat^+$ are well-defined $\rho$.
    %
    By induction hypothesis on patterns, the step function of each patterns $pat$ matching $v$
    produces a boolean $b'$ and an environment $\rho^{\prime\prime}$ where identifiers defined by
    $pat$ are well-defined if $b'$ is \true.
    %
    For all identifier $x \in \defs{pat^+}$, \Cref{rule:semantics:grsync:causality:defs:pat} implies
    that there exists a pattern $pat$ such that $x \in \defs{pat}$.
    %
    Compatible states for $e_g$ are empty, \cf \Cref{eq:semantics:grsync:opsem:stateTy:pats}.
    \Cref{lem:semantics:grsync:theorems:productivity:expression} ensures that the step function of $e_g$ on
    the environment $\rho \ctxconcat \bigCtxconcat \rho^{\prime\prime+}$ produces a well-defined value
    $b_g$, that is a boolean because the pattern is well-typed.
    %
    The step function of \guarded{pat^+ = v^+}{e_g}, as defined by
    \nameref{rule:semantics:grsync:opsem:step:pats:bool}, produces a boolean
    $b = b_g \bigwedge b^{\prime+}$ and an environment
    $\rho' = \bigCtxconcat \rho^{\prime\prime+}$. Therefore, if $b$ is \true, all $b'$ are
    \true and every identifier $x \in \defs{pat^+}$ has a well-defined value in $\rho'$.

    \paragraph{Wildcards}
    Let $\rho$ be an environment. The step function of $\_ = v$, as defined by \nameref{rule:%
        semantics:grsync:opsem:step:pat:wild}, produces a boolean $b$ that is always \true, and
    an environment $\rho'$ that is empty. \Cref{rule:semantics:grsync:causality:defs:pat:wild}
    ensures that \defs{\_} is the empty set. Therefore all identifiers defined by $\_$ are
    well-defined in $\rho'$.

    \paragraph{Identifiers}
    Let $\rho$ be an environment. The step function of $x = v$, as defined by \nameref{rule:%
        semantics:grsync:opsem:step:pat:ident}, produces a boolean $b$ that is always \true,
    and an environment $\rho' = [v / x]$. \Cref{rule:semantics:grsync:causality:defs:pat:ident}
    ensures that $\defs{x} = \{ x \}$. Therefore all identifiers defined by $x$ are well-defined in
    $\rho'$.

    \paragraph{Constants}
    Let $\rho$ be an environment. The step function of the constant pattern matching $c = v$, as
    defined by \nameref{rule:semantics:grsync:opsem:step:pat:const}, produces a boolean $b$ that
    equals $c \kwf{==} v$, and an environment $\rho'$ that is empty.
    \Cref{rule:semantics:grsync:causality:defs:pat:const} ensures that \defs{c} is the empty set.
    Therefore all identifiers defined by $c$ are well-defined in $\rho'$.

    \setcounter{paragraph}{0}
\end{skippableproof}

\begin{lemma}[Productivity of \GRSync event patterns]
    \label{lem:semantics:grsync:theorems:productivity:epat}
    Let \guarded{epat^+}{e_g} be a well-typed \GRSync event pattern guarded by $e_g$, let $\rho$ be
    an environment such that all dependencies of $epat^+$ and $e_g$, except those defined by
    $epat^+$ (\ie all identifiers in $\deps{epat^+} \cup \deps{e_g} \backslash \defs{epat^+}$), have
    well-defined values in the environment $\rho$ (non-$\bot$).
    %
    Assuming that no runtime errors occur, for all state $s$ corresponding to $epat^+$, the step
    function of \guarded{epat^+}{e_g} in $\rho$ from $s$ produces a new compatible state $s'$, a
    boolean $b$, and an environment $\rho'$ where identifiers defined by $epat^+$ (\ie all
    identifiers in \defs{epat^+}) are well-defined if $b$ is \true.
\end{lemma}
\begin{skippableproof}
    We proceed by structural induction on event patterns.

    \paragraph{Guarded event patterns}
    Let \guarded{epat^+}{e_g} be a guarded tuple of event patterns, let $s$ be a compatible state
    for the pattern, let $\rho$ be an environment such that all dependencies of $epat^+$ and $e_g$,
    except those from the pattern $epat^+$, are well-defined $\rho$.
    %
    \Cref{eq:semantics:grsync:opsem:stateTy:epats} ensures there exists a list of states $s_{epat}^+$
    compatible with $epat^+$ such that $s = (s_{epat}^+)$.
    %
    \Cref{rule:semantics:grsync:causality:dependencies:deps:epat} ensures that all the dependencies
    of each event pattern $epat$ have well-defined values in $\rho$. By induction hypothesis on
    event patterns, the step function of each event patterns $epat$ produces a boolean $b'$, a new
    compatible state $s'_{epat}$ and an environment $\rho^{\prime\prime}$ where identifiers defined
    by $epat$ are well-defined if $b'$ is \true.
    %
    For all identifier $x \in \defs{epat^+}$, \Cref{rule:semantics:grsync:causality:defs:epat}
    implies that there exists an event pattern $epat$ such that $x \in \defs{epat}$.
    %
    Compatible states for $e_g$ are empty, \cf \Cref{eq:semantics:grsync:opsem:stateTy:epats}.
    \Cref{lem:semantics:grsync:theorems:productivity:expression} ensures that the step function of $e_g$ on
    the environment $\rho \ctxconcat \bigCtxconcat \rho^{\prime\prime+}$ produces the well-defined value
    $b_g$, that is a boolean because the event pattern is well-typed.
    %
    The step function of \guarded{epat^+}{e_g}, as defined by
    \nameref{rule:semantics:grsync:opsem:step:epats:bool}, produces a boolean
    $b = b_g \bigwedge b^{\prime+}$, a new compatible
    state $s' = (s_{epat}^{\prime+})$ an environment $\rho' = \bigCtxconcat \rho^{\prime\prime+}$.
    Therefore, if $b$ is \true, all $b'$ are \true and every identifier $x \in
        \defs{epat^+}$ has a well-defined value in $\rho'$.

    \paragraph{Event occurrences}
    Let \detect{x}{y} be an event detection, let $s$ be a compatible state for the pattern, let
    $\rho$ be an environment such that all dependencies of the event-pattern are well-defined $\rho$.
    \Cref{rule:semantics:grsync:causality:dependencies:deps:epat:event} ensures that the value of
    the event $y$ is well-defined in $\rho$.
    %
    \Cref{eq:semantics:grsync:opsem:init:epat:event} ensures $s$ is an empty state.
    %
    If $\rho[y] = \none$ (event $y$ is absent) then the step function of \detect{x}{y},
    following the rule \nameref{rule:semantics:grsync:opsem:step:epat:event:none} produces a boolean
    $b$ that is \false and an environment $\rho' = \bot_{(x)}$ that maps $x$ to $\bot$.
    %
    If $\rho[y] = \some{v}$ (event $y$ is present) then the step function of \detect{x}{y}
    follows the rule \nameref{rule:semantics:grsync:opsem:step:epat:event:some} and produces an empty
    state and a boolean $b$ that is \true along with an environment $\rho' = [v / x]$ in which
    $x$ is well-defined.
    %
    \Cref{rule:semantics:grsync:causality:defs:epat:event} ensures that \defs{\detect{x}{y}} equals
    $\{ x \}$. Therefore all identifiers defined by \detect{x}{y} are well-defined in $\rho'$.

    \paragraph{Rising edges}
    Let $e$ be a rising edge event pattern, let $s$ be a compatible state for the pattern, let
    $\rho$ be an environment such that all dependencies of the event-pattern are well-defined $\rho$,
    and let $s$ be a corresponding state.
    %
    \Cref{rule:semantics:grsync:causality:dependencies:deps:epat:edge} ensures that all dependencies
    of the expression $e$ are well-defined in the environment $\rho$, and
    \Cref{eq:semantics:grsync:opsem:init:epat:edge} ensures that there exist a well-defined
    value $v_m$ and a state $s_e$ corresponding to $e$ such that $s = (v_m, s_e)$.
    \Cref{lem:semantics:grsync:theorems:productivity:expression} ensures the step function of $e$ from $s_e$ produces a new compatible state
    $s_e'$ and the well-defined value $b'$, that is a boolean because the event pattern is
    well-typed. The step function of the rising edge $e$, as defined by
    \nameref{rule:semantics:grsync:opsem:step:epat:edge:bool} produces the boolean
    $b = !v_m \wedge b'$, the compatible state $s' = (b', s_e')$, and an environment $\rho'$ that is
    empty. \Cref{rule:semantics:grsync:causality:defs:epat:edge} ensures that \defs{e} is the empty
    set. Therefore all identifiers defined by the rising edge $e$ are well-defined in $\rho'$.

    \setcounter{paragraph}{0}
\end{skippableproof}

\begin{lemma}[Productivity of \GRSync equations]
    \label{lem:semantics:grsync:theorems:productivity:equations}
    Let $eq^+$ be a well-typed and causal \GRSync equation, let $\rho$ be an environment such that
    all dependencies of $eq^+$ (\ie \deps{eq^+}) have a well-defined value in $\rho$ (non-$\bot$).
    %
    Assuming that no runtime errors occur, for all state $s$ corresponding to $eq^+$, the step
    function of $eq^+$ in $\rho$ from $s$ produces a new compatible state $s'$, and an environment
    $\rho'$ where all identifiers defined by $eq^+$ (\ie \defs{eq^+}) are well-defined in $\rho'$.
\end{lemma}
\begin{remark}
    In synchronous languages with explicit clocks, both clock calculus and causality analysis must
    be satisfied for a set of equations to ensure that it is productive. In \GRSync, the constructs
    \kwf{match} and \kwf{when} filter equations using pattern matching or event occurrences,
    but maintain the global equation on the base clock. As a result, only causality needs to be
    checked to establish the productivity of \GRSync equations.
\end{remark}
\begin{skippableproof}
    The proof proceeds by structural induction on equations.

    \paragraph{Overview of the proof strategy}
    The difficult points of the proof are the following.
    \begin{itemize}
        \item For \kwf{match} equations, the exhaustiveness of patterns (verified by the \Rust
              compiler on the generated code) ensures that a branch will always match.
        \item For \kwf{when} equations, signals not updated in a branch retain their last value, and
              events not emitted are set to \none, ensuring an available value for every identifier.
        \item For blocks of mutually recursive equations, we use rank induction
              (\Cref{def:semantics:grsync:causality:induction:rank}) to prove that
              identifiers are iteratively defined following the rank (\Cref{def:semantics:grsync:causality:rank}).
              The rank is the maximum length of the dependency sequence of an identifier in a causal
              component.
    \end{itemize}

    \paragraph{Declarations}
    Let $x^+ = e$ be a well-typed declaration, let $s$ be a compatible state for the equation, let
    $\rho$ be an environment such that all dependencies of the equation are well-defined in $\rho$.
    %
    \Cref{rule:semantics:grsync:causality:dependencies:deps:eq:decl} ensures that all dependencies
    of $e$ are well-defined in $\rho$, and \Cref{eq:semantics:grsync:opsem:init:eq:decl} ensures
    that the state $s$ corresponds to $e$.
    %
    \Cref{lem:semantics:grsync:theorems:productivity:expression} ensures the step function of $e$ from $s$
    produces the well-defined values $v^+$. The step function of $x^+ = e$, following the rule
    \nameref{rule:semantics:grsync:opsem:step:eq:decl} produces an environment $\rho' =
        [(v / x)^+]$, in which the values of $x^+$ are well-defined. \Cref{rule:semantics:%
        grsync:causality:defs:eq:decl} ensures that \defs{x^+ = e} equals $\{ x^+ \}$. Therefore all
    identifiers defined by $x^+ = e$ are well-defined in $\rho'$.

    \paragraph{Match}
    Let \match{e}{x^+}{(\branch{pat_i^+}{e_i}{eq_i^+})^+} be a well-typed \kwf{match} equation, let
    $s$ be a compatible state for the equation, let $\rho$ be an environment such that all
    dependencies of the equation are well-defined in $\rho$.
    %
    \Cref{rule:semantics:grsync:causality:dependencies:deps:eq:match} ensures that all dependencies
    of $e$, $e_i^+$, and $(eq_i^+)^+$, except those defined by $pat_i^+$, are well-defined in
    $\rho$, and \Cref{eq:semantics:grsync:opsem:stateTy:eq:match} ensures that there exist a state
    $s_e$ corresponding to $e$ and a tuple of states $(s_i^+)$ each corresponding to the blocks of
    equations $eq_i^+$ such that $s = (s_e, s_i^+)$.
    %
    \Cref{lem:semantics:grsync:theorems:productivity:expression} ensures the step function of $e$ from $s_e$
    produces the well-defined values $v_e^+$.
    %
    \Cref{lem:semantics:grsync:theorems:productivity:patterns} ensures the step function of
    \guarded{pat_i^+ = v_e^+}{e_i} in $\rho$ from an empty state produces another
    empty state, a boolean $b_i$, and an environment $\rho_i$ where identifiers defined by $pat_i^+$
    are well-defined if $b_i$ is \true.
    %
    Because patterns are exhaustive, \emph{there exists $k$ such that $b_k$ is \true} and for all
    $i < k$, $b_i$ is \false. Consequently, $\rho_k$ defines the identifiers defined by
    $pat_k^+$.
    %
    Therefore $\rho \ctxconcat \rho_k$ defines all dependencies of $eq_k^+$, so by induction hypothesis%
    , its step function from state $s_k$ produces a new state $s'_k$ and an environment $\rho'$
    where identifiers defined in $eq_k^+$ are well-defined.
    %
    The rule \nameref{rule:semantics:grsync:opsem:step:eq:match:value} states that
    the step function of the \kwf{match} equation returns the environment $\rho'$. Because the
    \kwf{match} equation is well-typed,
    the rule \nameref{rule:grust:id_usage:grsync:eq:declared:match} on identifier usage ensures that
    for all equations-branch $eq_i^+$ we have
    $\defs{eq_i^+} = \{x^+\}$. Therefore all identifiers defined by the \kwf{match} equation are
    well-defined in $\rho'$.

    \paragraph{When}
    Let \when{x^+}{x_m^+ = c^+}{(\branch{epat_i^+}{e_i}{eq_i^+})^+} be well-typed, let $s$ be a
    compatible state for the equation, let $\rho$ be an environment such that all dependencies of
    the equation are well-defined in $\rho$.
    %
    \Cref{rule:semantics:grsync:causality:dependencies:deps:eq:when} ensures that all dependencies
    of $(epat_i^+)^+$, $e_i^+$, and $(eq_i^+)^+$, except those defined by $epat_i^+$ and the
    memories $(\last{x_m})^+$, are well-defined in $\rho$, and \Cref{eq:semantics:grsync:opsem:init%
        :eq:when} ensures that there exist a tuple of well-defined values $(v_m^+)$, a tuple of
    states $(p_i^+)$ each corresponding to $epat_i^+$, and a tuple of states $(s_i^+)$ each
    corresponding to the blocks of equations $eq_i^+$ such that $s = (v_m^+, p_i^+, s_i^+)$.
    %
    \Cref{lem:semantics:grsync:theorems:productivity:epat} ensures the step function of
    \guarded{epat_i^+}{e_i} in $\rho$ from $p_i$ produces a new state $p_i'$, a boolean $b_i$, and
    an environment $\rho_i$ where identifiers defined by $epat_i^+$ are well-defined if $b_i$ is
    \true.

    Suppose there exists $k$ such that $b_k$ is \true and for all $i < k$, $b_i$ is
    \false. Consequently, $\rho_k$ defines the identifiers defined by $epat_k^+$.
    %
    The environment $\rho \ctxconcat \rho_k \ctxconcat \rho_m \ctxconcat \rho_k'$ defines all dependencies of
    $eq_k^+$, with $\rho_m$ being the memory of signals defined by the \kwf{when} equation
    ($[(v_m/\last{x_m})^+]$) and $\rho_k'$ being the environment of identifiers not defined by
    $eq_k^+$ (signals get their \kwf{last} value and events are set to \none).
    %
    By induction hypothesis, the step function of $eq_k^+$ from state $s_k$ produces a new state
    $s'_k$ and an environment $\rho'$ where identifiers defined in $eq_k^+$ are well-defined.
    %
    In that case, the rule \nameref{rule:semantics:grsync:opsem:step:eq:when:branch} is applied, and
    the step function of the \kwf{when} equation returns the environment $\rho_m \ctxconcat \rho'_k
        \ctxconcat \rho'$, where the identifiers defined by the \kwf{when} equation are well-defined
    (\cf \Cref{rule:semantics:grsync:causality:defs:eq:when}).

    In the case all $b_i^+$ are \false, the step function of the \kwf{when} equation follows
    the rule \nameref{rule:semantics:grsync:opsem:step:eq:when:default} and returns the environment
    $\rho'$ which gives default values to all identifiers defined by the \kwf{when} equation:
    signals get their last value and events are set to \none.

    \paragraph{Mutually recursive equations}

    Let $n>0$, let $eq^+$ be a well-typed and causal set of $n$ equations, let $s$ be a compatible
    state for the set of equations, and let $\rho$ be an environment such that all dependencies of
    the equation are well-defined in $\rho$ (\ie for all external dependency $y \in \deps{eq^+}$ we
    have $\rho[y] \in \Val$).

    To proceed, we rely on the following definitions.
    Let \Val denote the set of all possible values and \ValBot the set of bottom values.
    Let $P_{\{x^+\}} = \prod_x \{x\}\times\ValBot$ be the set of environments over the finite set of
    identifiers $\{x^+\}$.
    The bottom environment $\botrho{\{x^+\}} \in P_{\{x^+\}}$ assigns $\bot$ to every identifier
    in $\{x^+\}$.

    \Cref{eq:semantics:grsync:opsem:stateTy:eqs} ensures that there exist a tuple of states
    $(s^{\prime+})$, each corresponding to one equation of $eq^+$, such that $s = (s^{\prime+})$.

    The step function of $eq^+$ uses the fix-point function that performs $n+1$ iterative step
    transitions of each individual equation, as stated by the rule
    \nameref{rule:semantics:grsync:opsem:step:eqs}.
    %
    From the state $s = (s^{\prime+})$ and an environment $\rho \in P_{(\deps{eq^+})}$, the fix-point
    function produces the sequence of $n+1$ environments $(\rho_k^+)$ in $P_{\defs{eq^+}}$, and a
    sequence of $n+1$ states $(s_k^{\prime+})$ defined as follows:
    \begin{equation*}
        \fixsem{k}{\rho}{s^{\prime+}}{eq^+}{s_k^{\prime+}}{\rho_k}
    \end{equation*}
    equivalently defined as follows:
    \begin{align*}
        \rho_0 & = \botrho{\defs{eq^+}} & \rho_{k+1} & = \bigCtxconcat (\rho'_{k+1})^+
    \end{align*}
    where each $\rho'_{k+1}$ is obtained by evaluating each $eq \in \{eq^+\}$ under the environment
    $\rho \ctxconcat \rho_k$:
    \begin{equation*}
        \eqOpsem{\rho \ctxconcat \rho_k}{s'}{eq}{s'_k}{\rho'_{k+1}}.
    \end{equation*}
    Intuitively, $\rho_k$ captures value of the identifiers defined by $eq^+$ at the $k^\text{th}$
    iteration. Beginning with the bottom empty environment where all identifiers are undefined (\ie
    equal to $\bot$), it is incrementally extended as more identifiers are determined.

    We want to prove that for all $x \in \defs{eq^+} \cup \deps{eq^+}$ the property \P{x} holds:
    \begin{equation*}
        \P{x}: \forall r \in \{0, \dots, n\}, \, \rank{eq^+}{x} \le r \Rightarrow (\rho \ctxconcat \rho_{r+1})[x] \in \Val.
    \end{equation*}

    We reason by rank induction (\Cref{def:semantics:grsync:causality:induction:rank}): we prove the
    property incrementally for identifiers following their rank.
    \begin{itemize}
        \item \textit{Base case:} Let $x \in \defs{eq^+} \cup \deps{eq^+}$ such that
              $\rank{eq^+}{x} = 0$. \Cref{def:semantics:grsync:causality:rank}
              gives the formal definition of the rank, it ensures that $x$ has no dependency in
              $eq^+$ (because the set
              $D_x = \{y \in \defs{eq^+} \cup \deps{eq^+} \mid \depOrd{eq^+}{x}{y}\}$ is empty).

              Let $r \in \{0, \dots, n\}$.
              We want to prove that, because it has no dependencies, $x$ is defined in $\rho \ctxconcat \rho_{r+1}$
              (\ie $(\rho \ctxconcat \rho_{r+1})[x] \in \Val$).

              If $x \notin \defs{eq^+}$ then $x$ is in \deps{eq^+}, which ensures that
              $\rho[x] \in \Val$, and thus $(\rho \ctxconcat \rho_{r+1})[x] \in \Val$.

              If $x \in \defs{eq^+}$, we want to prove that the equation defining $x$ has no dependencies.
              Let us examine the equation $eq_x$ that defines $x$ (that exists thanks ot
              \Cref{rule:semantics:grsync:causality:defs:eq}). We want to prove that \deps{eq_x} is
              an empty set.
              \begin{itemize}
                  \item Suppose $eq_x$ is a declaration $z^+ = e$.
                        \Cref{rule:semantics:grsync:causality:defs:eq:decl} ensures $x$ is in
                        $\{z^+\}$.
                        \Cref{rule:semantics:grsync:causality:dependencies:rel:eq:decl} states that
                        $D_x = \deps{e}$, which implies that \deps{e} is empty.
                        \Cref{rule:semantics:grsync:causality:dependencies:deps:eq:decl} ensures that
                        \deps{eq_x} is empty.
                  \item Suppose $eq_x$ is a \kwf{match} equation:
                        $$\match{e}{z^+}{(\branch{pat_i^+}{e_g}{eq_i^+})^+}.$$
                        \Cref{rule:semantics:grsync:causality:defs:eq:match} ensures $x$ is in
                        $\{z^+\}$. \Cref{rule:semantics:grsync:causality:dependencies:rel:eq:match}
                        states that $D_x = \deps{e} \cup \bigcup (\deps{e_i} \cup \deps{eq_i^+})^+$
                        (the renaming functions $\sigma^{x^+}_{pat_i^+}$ induce relation bridges between
                        the $z^+$ and the dependencies from the branch-equations $eq_i^+$), includes
                        the dependencies of the \kwf{match} equation (\Cref{rule:semantics:grsync:%
                            causality:dependencies:deps:eq:match}), which is then empty.
                  \item Suppose $eq_x$ is a \kwf{when} equation:
                        $$\when{z^+}{x_m^+ = c^+}{(\branch{epat_i^+}{e_g}{eq_i^+})^+}.$$
                        \Cref{rule:semantics:grsync:causality:defs:eq:when} ensures $x$ is in
                        $\{z^+, \last{x_m}^+\}$. \Cref{rule:semantics:grsync:causality:dependencies%
                            :rel:eq:when} states that $D_x$ equals $\bigcup (\deps{epat_i} \cup
                            \deps{e_i} \cup \deps{eq_i^+})^+$ (the renaming functions
                        $\sigma^{x^+}_{epat_i^+}$ induce relation bridges between the $z^+$ and the
                        dependencies from the branch-equations $eq_i^+$), includes the dependencies
                        of the \kwf{when} equation (\Cref{rule:semantics:grsync:causality:%
                            dependencies:deps:eq:when}), which is then empty.
              \end{itemize}
              Therefore, \deps{eq_x} is an empty set.

              The structural induction hypothesis on equations gives:
              \begin{equation*}
                  \eqOpsem{\rho \ctxconcat \rho_r}{s_x}{eq_x}{s'_x}{\rho'_{r+1}},
              \end{equation*}
              with $\rho'_{r+1}[x] \in \Val$.
              %
              Because $(\rho \ctxconcat \rho_{r+1})[x]$ equals $\rho_{r+1}[x]$ which equal
              $(\bigCtxconcat (\rho'_{r+1})^+)[x]$ and finally equals $\rho'_{r+1}[x]$, we can
              state that $(\rho \ctxconcat \rho_{r+1})[x] \in \Val$ and that \P{x} holds.

        \item \textit{Inductive step:} Let $k \in \N$.
              We assume \P{y} for all $y \in \defs{eq^+} \cup \deps{eq^+}$ such that $\rank{eq^+}{y} \le k$.
              Let $x \in \defs{eq^+} \cup \deps{eq^+}$ such that $\rank{eq^+}{x} = k+1$.
              Let $r \in \{0, \dots, n\}$ such that $k+1 \le r$.

              For all $y \in \defs{eq^+} \cup \deps{eq^+}$ such that $x$ depends on $y$ (\ie \depOrd{eq^+}{x}{y}),
              \Cref{lem:semantics:grsync:causality:rank:decrease}
              states that $\rank{eq^+}{x} > \rank{eq^+}{y}$, which implies that $\rank{eq^+}{y} \le k$.
              By rank induction hypothesis, for all $y \in \defs{eq^+} \cup \deps{eq^+}$ such that
              \depOrd{eq^+}{x}{y}, \P{y} holds (\ie for all $r' \in \{0, \dots, n\}$ such that
              $\rank{eq^+}{y} \le r'$ we have $(\rho \ctxconcat \rho_{r'+1})[y] \in \Val$).
              Because $\rank{eq^+}{y} \le k \le r-1$, then for all $y \in \defs{eq^+} \cup \deps{eq^+}$ such that
              \depOrd{eq^+}{x}{y} we have $(\rho \ctxconcat \rho_r)[y] \in \Val$.

              The equation $eq_x$ that defines $x$ only depends on identifiers that are well-defined
              in $\rho \ctxconcat \rho_r$. The induction hypothesis on equations gives:
              \begin{equation*}
                  \eqOpsem{\rho \ctxconcat \rho_r}{s_x}{eq_x}{s'_x}{\rho'_{r+1}},
              \end{equation*}
              with  $\rho'_{r+1}[x] \in \Val$.
              %
              Because $(\rho \ctxconcat \rho_{r+1})[x]$ equals $\rho_{r+1}[x]$ which equal
              $(\bigCtxconcat (\rho'_{r+1})^+)[x]$ and finally equals $\rho'_{r+1}[x]$, we can
              state that $(\rho \ctxconcat \rho_{r+1})[x] \in \Val$ and that \P{x} holds.
    \end{itemize}

    \Cref{lem:semantics:grsync:causality:rank:bounded} ensures that for all $x \in \defs{eq^+}$,
    we have $\rank{eq^+}{x} \le n$, where $n$ is the number of equations. Then, all
    $x \in \defs{eq^+}$ ensures $\rho_{n+1}[x] \in \Val$,
    \ie all identifiers of $f$ are defined after the $(n+1)^\text{th}$ step iteration, in particular
    all $x_m$ and $x_o$.

    \paragraph{Other approaches}
    We give a brief overview of how other synchronous languages address the productivity.

    \begin{description}
        \item[Fixed-point approach]
              The operational semantics of a set of equations is usually represented as a fix-point
              application of the step functions of each individual equation.
              A common method for proving its productivity involves applying Kleene fixed-point theorem, as
              shown in prior work~\cite{DBLP:journals/tecs/ColacoMPP23}.
              The causality of the set of equations is used to ensure a fix-point where all identifier are
              non-bottom.

        \item[Causal induction]
              Causal induction~\cite{bourke:hal-03936656} is another proof technique used to reason about
              properties of mutually recursive equations by following the causality (dependency) order
              directly. This approach is very similar to rank induction.

        \item[Fixed scheduling]
              Another approach is to use the causality of the set of equations to derive a scheduling,
              an order in which equations can be called in the operational semantics to produce their values.
              This method is discussed in \Cref{rmk:semantics:grsync:opsem:step:eqs}.
              We want the operational semantics to describe the execution of the high-level program,
              while a scheduling would describe the execution of the compiled code.
    \end{description}

    \setcounter{paragraph}{0}
\end{skippableproof}

\begin{theorem}[Productivity]
    \label{thm:semantics:grsync:theorems:productivity}
    Let $f$ be a well-typed and causal component in a program \G with a typing context \tyCtx, and
    assuming no mutually recursive components exist. Assuming that no runtime errors occur, for all
    input lists $L^+$ of length $n \in \Nstar$, \obsComp{f}{L^+}{n} is well-defined.
\end{theorem}
\begin{skippableproof}
    Let \component{f}{x^+_i}{x^+_o}{eq^+}{x^+_m = c^+} be a well-typed and causal component in a
    program \G and a typing context \tyCtx, let $(v_m^+, s^+)$ be a compatible state (defined
    in \Cref{eq:semantics:grsync:opsem:stateTy:comp}), and let $v_i^+$ be input values.

    Because $f$ is well-typed, the dependencies of its internal equations $eq^+$ are included in
    the set of inputs and memories $\{x_i^+, \last{x_m}^+\}$. Therefore, the causality of $f$ allows
    to apply \Cref{lem:semantics:grsync:theorems:productivity:equations}, ensuring that the step function of
    $eq^+$, from a state $(s^+)$ and a context $[(v_i/x_i)^+, (v_m/\last{x_m})^+]$ produces a new
    compatible state $(s^{\prime+})$ and an environment $\rho$ where all identifiers from
    \defs{eq^+} are well-defined.

    And because $f$ is well-typed, all $x_o^+$ and $x_m^+$ are in \defs{eq^+}, and thus are
    well-defined in $\rho$. The rule \nameref{rule:semantics:grsync:opsem:step:comp} then states
    that the output values of the step function of $f$ are well-defined and the next state is
    compatible.

    \setcounter{paragraph}{0}
\end{skippableproof}

\paragraph{Summary}
In summary, this section established the productivity of well-formed \GRSync components. The central
result, presented in \Cref{thm:semantics:grsync:theorems:productivity}, guarantees that any
well-typed and causal component produces a well-defined output for any sequence of inputs, assuming
no runtime errors.
%
This property was demonstrated through a proof by structural induction on the language constructs.
%
The key challenge, posed by mutually recursive equations, was addressed using rank induction
(\Cref{def:semantics:grsync:causality:induction:rank}), which ensures that all identifiers are
computed in a finite number of steps by following dataflow dependencies.
%
This formal guarantee of productivity is a cornerstone of the language design, ensuring that every
component behaves as a total function over input streams and will not deadlock or halt unexpectedly
during its execution.
